\section{Research question, source audit, and precise contribution}
\label{sec:scope}

The nonlinear feedback equation
\begin{equation}\label{eq:model}
 uv(r)-\frac{v(r)^2}{2}=\Li_s\bigl(r\ee^{v(r)}\bigr),
 \qquad v(0)=0,\qquad s>0,\quad u>0,
\end{equation}
provides a concrete meeting point for polylogarithms and transseries.
The polylogarithm supplies a boundary singularity whose exponent changes
with $s$.  Nonlinear inversion may reach that singularity before, after,
or exactly when the ordinary inverse derivative vanishes.  The equality
case produces fractional powers, logarithms, and iterated logarithms;
it is especially well suited to a precise distinction between a formal
expansion and a convergent analytic chart.

The source report \cite{RepoFeedback}, Theorem \texttt{thm:feedback},
proves the case in which an interior characteristic point exists.
Its first question in Section \texttt{sec:future} asks for the remaining
boundary cases and prints the supercritical coefficient constant as an
unproved candidate.  We prove that candidate and all critical cases.
This is a resolution of an explicitly stated project question, not a
claim to settle a previously established open problem in the wider
literature.

\subsection{Snapshot and comparison sources}

The repository was inspected at commit
\begin{center}
\small\texttt{6207c3d40047fd8d1400a4d509fee71c9a3a8132}.
\end{center}
The review covered the two topic indexes, the incoming-report index,
the five incoming ZIP articles, and the relevant critical-feedback
and endpoint articles in the Transseries collection.  It was a focused
mathematical audit, not a line-by-line verification of the complete
552-page Polylogarithms manuscript or of every Transseries report.
The source inventory and hashes are recorded in \texttt{provenance.json}.

\begin{longtable}{@{}>{\raggedright\arraybackslash}p{0.25\textwidth}>{\raggedright\arraybackslash}p{0.36\textwidth}>{\raggedright\arraybackslash}p{0.33\textwidth}@{}}
\toprule
Source & Result already available & Role in this report\\
\midrule
\endhead
Incoming feedback report \cite{RepoFeedback}
& Positive finite coefficients, exact interior radius and fold,
and the boundary conjecture
& Interior theorem credited; boundary conjecture proved and extended\\
Critical Hahn report \cite{RepoHahn}
& Linear feedback with power tails, critical fractional inverses,
stable coefficient laws, coupling crossover
& General mechanism and comparison model\\
Gaussian endpoint report \cite{RepoGaussian}
& Uniform order crossover for eventually exact power tails,
Lambert coordinates and local limits
& Antecedent for the order--logarithm scaling\\
Other four incoming articles \cite{RepoJets,RepoUniform,RepoNormalized,RepoFolds}
& Resonant order jets, ordinary polylogarithm inverse charts,
complex critical values and folds
& Prevents presenting those existing charts as new results\\
\bottomrule
\end{longtable}
\addtocounter{table}{-1}% This uncaptioned comparison is not a numbered table.

The nonlinear function arising here is
\begin{equation}\label{eq:fdef-intro}
 f_u(t)=u-\sqrt{u^2-2\Li_s(t)}.
\end{equation}
Even when its coefficients have a power-law tail, they are not in
general eventually an exact multiple of $n^{-s}$.  Thus the existing
theorems for a polylogarithm plus a finite polynomial do not by
themselves verify this model's endpoint continuation or its correction
constants.  We give those proofs explicitly.

\subsection{Classical inputs and novelty boundary}

We use the classical convergent Jonqui\`ere expansion
\begin{equation}\label{eq:jonquiere}
 \Li_s(\ee^{-x})=
 \Gamma(1-s)x^{s-1}+
 \sum_{k\ge0}\zeta(s-k)\frac{(-x)^k}{k!},
 \qquad 0<|x|<2\pi,
\end{equation}
for nonintegral positive $s$, with a fixed branch of $\Log x$.
At positive integers the colliding Gamma and zeta terms are combined
before the parameter limit is taken.  These are classical identities,
recorded in DLMF \S25.12 \cite{DLMFPolylog}.
All displayed local series below are restricted to sufficiently small
sectors within this convergence domain.

The conversion of local singular expansions into Taylor coefficient
asymptotics is the singularity-analysis theory of Flajolet and Odlyzko
\cite{FO}.  The relation between implicit generating functions,
conditioned sums, and stable laws is also classical; Janson
\cite{Janson}, especially Sections 17--18, provides an extensive
treatment.  We include the particular Fourier proof needed for a
varying feedback parameter, so that the uniformity is not left to an
unstated invocation of a fixed-distribution theorem.

The exact boundary formula, the nonlinear critical charts, the
order-three correction polynomials, and the explicit uniform order
chart are the results established here for \eqref{eq:model}.
Their novelty relative to the inspected repository is documented above.
A focused literature search found the standard general mechanisms and
their antecedents; it did not establish publication priority for every
model-specific formula.  No theorem is advertised as a new general
principle of probability or transseries.

\subsection{Conventions}

Write
\[
 a_n(s,u)=\coef r n v(r),\qquad
 A_s=\zeta(s-1)\quad(s>2),\qquad
 \Uc(s)=\sqrt{2\zeta(s)+A_s^2}.
\]
An assertion about the coefficients $a_n$ always concerns the germ
$v(0)=0$ of \eqref{eq:model}.  It is not automatically an assertion
about the ordinary compositional inverse of $-\Li_s(-\ee^x)$.
The incoming source already makes this distinction; the special
dilogarithm identification does not extend to arbitrary $s$.

All square roots on a positive real argument denote the positive root.
Local inverse branches are selected by positivity for positive real
variables and then continued to a slit sector.  A \emph{convergent
transseries chart} means a holomorphic function of finitely many small
coordinates whose ordinary multivariable Taylor series converges
absolutely after substitution.  This claim is stronger than a
Poincar\'e asymptotic expansion.  The asymptotic coefficient expansions
obtained after singularity transfer need not themselves converge.
